\runningheader{Computer Science Principles}{AP Practice Exam (Data)}{Page \thepage\ of \numpages}

\begingradingrange{cspdata}

\question[5]
Computers process instructions iteratively. Given:
\begin{cspcode}
a <- 38
b <- 54
a <- b + a
DISPLAY(a + b)
\end{cspcode}
How are these lines processed?
\begin{choices}
  \choice All lines at once
  \choice Metadata first
  \CorrectChoice One line at a time, in order
  \choice Simultaneous display while processing
\end{choices}
\begin{solution}
  Sequential execution processes statements in order.
\end{solution}

\question[5]
A whale collar records only timestamped location. Which cannot be answered from that data alone?
\begin{choices}
  \choice Weekly distance traveled
  \choice Typical region visited
  \choice Whether nearby tracked whales moved together
  \CorrectChoice Whether behavior changes with weather
\end{choices}
\begin{solution}
  Weather requires an additional data source.
\end{solution}

\question[5]
Which is predictive use of large data?
\begin{choices}
  \choice Monthly billing totals
  \CorrectChoice Product recommendations from past purchases
  \choice Largest spending addresses
  \choice Most active purchase hour
\end{choices}
\begin{solution}
  Recommendations use prior behavior to predict future choices.
\end{solution}

\question[5]
Which is \textbf{not} metadata for an SMS message?
\begin{choices}
  \choice Timestamp
  \choice Sender/receiver phone numbers
  \choice Sender/receiver location
  \CorrectChoice Message text content
\end{choices}
\begin{solution}
  Content is primary data, not metadata.
\end{solution}

\newpage 

\question[5]
Google Trends data is most useful for:
\begin{choices}
  \CorrectChoice Detecting rising public interest in weather events
  \choice Determining cheapest hotel booking day
  \choice Pricing a chainsaw
  \choice Determining Iditarod length
\end{choices}
\begin{solution}
  Search-frequency trends reveal shifts in public interest.
\end{solution}

\question[5]
Which cannot be answered from student-interest fields (GPA, intended major, target colleges, school, address) alone?
\begin{choices}
  \CorrectChoice Number of admitted students
  \choice Number interested in each major
  \choice Interest by high school
  \choice Average GPA by target-college group
\end{choices}
\begin{solution}
  Admissions outcomes are not present.
\end{solution}


\question[5]
Why can larger search engines produce more useful trend data? \textbf{Select two answers.}
\begin{checkboxes}
  \CorrectChoice Larger sample size improves signal quality
  \CorrectChoice Broader demographics reduce sampling bias
  \choice Larger size guarantees perfect accuracy
  \choice Popularity-in-results always guarantees correctness
\end{checkboxes}
\begin{solution}
  Sample size and representativeness are key strengths.
\end{solution}

\question[5]
Which search-engine capability is most unique?
\begin{choices}
  \choice Listing creator names
  \CorrectChoice Filtering across many sources quickly
  \choice Finding print publication dates
  \choice Returning one canonical result only
\end{choices}
\begin{solution}
  Aggregated indexing and filtering across sites is central.
\end{solution}

\newpage 

\question[5]
In a library catalog, which filter is least useful?
\begin{choices}
  \choice Genre
  \choice Branch availability
  \choice Author
  \CorrectChoice Number of letters in title
\end{choices}
\begin{solution}
  Title length is rarely meaningful for discovery.
\end{solution}

\question[5]
For a daily business spreadsheet, which two features matter most? \textbf{Select two answers.}
\begin{checkboxes}
  \CorrectChoice Graph generation
  \CorrectChoice Auto-updating totals/formulas
  \choice Map charts
  \choice Advanced formula linting only
\end{checkboxes}
\begin{solution}
  Visualization and automatic recalculation are high-value basics.
\end{solution}

\question[5]
Which best improves transmission efficiency of large data?
\begin{choices}
  \choice Encryption
  \CorrectChoice Compression
  \choice Emulation
  \choice Keep-alive pings
\end{choices}
\begin{solution}
  Compression reduces payload size.
\end{solution}

\question[5]
Which is the standard way to protect sensitive personal data?
\begin{choices}
  \CorrectChoice Encryption
  \choice Compression
  \choice Emulation
  \choice Keep-alive pings
\end{choices}
\begin{solution}
  Encryption protects confidentiality in transit and at rest.
\end{solution}

\endgradingrange{cspdata}
